\pagebreak

{
\thispagestyle{empty}

\mbox{}

\vspace{4cm}

\footnotesize
\noindent{}\emph{\noindent{}O ato de escrever, a incisão feita ao gravar letras, palavras e frases sobre uma página em branco, cria mundos, realidades e universos por completo. Este incansável e interminável movimento é feito neste exato instante e neste exato documento\footnote{Pretende-se que as diversas interpretações dos textos, aqui expostos, ressoem em você que lê, o leitor.}. E assim, como intérprete de uma longa tradição de textos e de livros, acredito que o nosso mundo --- e todas as coisas que ele contém --- está sendo criado e reinventado por nós a todo instante. Por fim, os textos aqui escritos e inscritos --- suas notas de rodapé\footnote{A disposição gráfica das notas de rodapé alude à página dos textos judaicos religiosos, bem como do pensamento rabínico que há no Talmud. Pois, no Talmud, a menção de um nome, situação ou ideia permite a introdução de uma história que alivia o clima de discussões complexas e leva o debate adiante. As notas, ou melhor, os comentários e as interpretações rabínicas, rondam, circunscrevem e flutuam ao redor do texto principal sem hierarquização. Aludindo à organização talmúdica, procuro incitar os outros a refletir, abrindo espaço para as mais diversas interpretações, histórias e vozes.} e suas interpretações (ressonâncias) que levam a outras histórias, anedotas e contextos --- atestam a veracidade da fala de Rabi Tarfon: ``Não lhe é exigido que complete a tarefa, mas não é livre para escapar dela''\footnote{\emph{Pirkei Avot} 2:16, escrito cerca do século \textsc{i}.}. A diagramação destas páginas-territórios são releituras dos textos sagrados, uma tentativa de apropriação das chispas e das fagulhas que voam e contornam o território do texto --- já que o espaço do texto é um território. E mais que nada, neste instante, é hora de nos apropriarmos dele.
}}

\pagebreak

\chapter{Novo capítulo}

\noindent{}Para a escritora e poeta chicana Gloria Anzaldúa (1942--2004) o termo
\emph{nepantla} diz respeito ao momento da escrita em que o todo é
caótico, ou melhor:

\begin{quote}
é um dos estágios da escrita, o estágio em que você tem todas essas
ideias, todas essas imagens, frases e parágrafos, e onde você tenta
transformá-los em uma única peça, uma história, enredo ou o que quer que
seja\footnote{\textsc{anzaldúa}, \emph{Borderlands/\,La Frontera}, p.\,237. Tradução da autora.}.
\end{quote}

\emph{Nepantla} é um espaço-tempo da construção da escrita, um termo que
se refere ao processo de estar entre meios. É um limiar em que múltiplas
formas de realidades (e de escritas) são vistas ao mesmo tempo ---
possivelmente enquanto um longínquo e interminável ensaio. Pois
``um ensaio é uma tentativa de incitar os outros a refletirem, de
levá-los a escrever complementos. (\ldots{}) Ele deve fazer rolar uma bola de
neve, na qual os complementos encubram cada vez mais a exposição
original''\footnote{\textsc{costa} \emph{apud} \textsc{flusser}, \emph{op.\,cit.}, p.\,11.}.

Segundo Gloria Anzaldúa\footnote{\textsc{anzaldúa}, \emph{op.\,cit.}, p.\,237. Tradução da autora.}, \emph{nepantla} é uma palavra do universo indígena
\emph{nahuatl} e indica algo que está entre dois corpos d'água --- um
espaço-tempo entre dois mundos. Este território é um limiar, um local
que não está nem aqui, nem lá, mas em transição.

Se as coisas são escritas entre meios, então \emph{nepantla} é um
espaço-tempo de fronteira, em que continuamente (e constantemente) se
encontra o mundo. São espaços-territórios --- geográficos ou
linguísticos --- de encontro com a alteridade. Assim, habitar o entre
resulta numa visão dupla, primeiro da perspectiva de uma cultura, depois
da perspectiva de outra. Ao viver entre fronteiras, as comunidades
judaicas, por exemplo, bem como os demais grupos socialmente
minoritários (chicanos, quilombolas, indígenas, refugiados etc.), acabam
vendo duplamente, interagindo com estes ambientes tanto no âmbito do
texto, quanto no âmbito do espaço, e criando locais de pertencimento.

\begin{center}\adforn{49}\end{center}

\noindent{}Na tradição da escrita judaica religiosa --- e de suas interpretações
---, as páginas, os papéis e os pergaminhos --- bem como os seus textos
--- são territorialidades de apropriação e espaços sagrados. Este texto
escrito e inscrito sobre um pedaço de papel, pano, couro ou celulose
influencia as comunidades, suas relações com a divindade e com os
mandamentos divinos. Nos mundos das narrativas e dos livros, a
manifestação do texto no espaço\footnote{As mais diversas interpretações
  dos textos sagrados judaicos se manifestam no espaço, pois conduzem as
  comunidades judaicas na diáspora e fora dela e acarretam mudanças
  reais nas maneiras de manifestar a judeidade, bem como nas dinâmicas
  dos territórios em que as comunidades se agregam.} --- como é o caso
da leitura dos rituais e das festividades judaicas religiosas ---
transforma as comunidades, seus espaços de agregação, bem como as suas
histórias. Em uma perspectiva de fronteira linguística, a ficção dos
textos, segundo o pesquisador e filósofo Leonel Olímpio:

\begin{quote}
Também implica no problema de como aquele mundo ficcional nos
influencia, e também até que ponto nós influenciamos aquele mundo
ficcional. Como assim? Por exemplo, se partirmos do pressuposto de que o
mundo daqui (presumidamente ``real'') não influencia o outro, toda
leitura, toda linguagem exposta em qualquer texto não mudaria de
interpretação com o tempo\footnote{OLIMPIO, \emph{Há o Anti-Limite?
  Sobre ficções, equívocos e metafísica experimental}, p.\,67.}.
\end{quote}

Neste sentido, os textos judaicos sagrados são passíveis de
interpretações, e o leitor é um participante ativo das páginas --- o que
faz com que os escritos se tornem territórios fluidos, espaços de
diálogo e de apropriação. O leitor de textos judaicos não lê
simplesmente palavras escritas em uma página; o leitor é um ator ativo e
um participante vocal na conversa.

O debate das escrituras ocorre em meio às páginas, e seus assuntos
abrangem séculos, milênios e gerações de professores, mestres, rabinos e
pensadores talmúdicos. O Talmud (em português: ``aprendizado''),
livro composto pela \emph{Mishná} (``estudo repetido'') e pela
\emph{Guemará} (``conclusão''), é uma série de textos-ensaios que
procuram --- a todo custo --- desmembrar as histórias (\emph{reductio ad
absurdum}), as regras, as restrições e as leis judaicas religiosas a
ponto de dilacerar suas minúcias.

Os textos sagrados, bem como as interpretações rabínicas, são escritos
às vezes em hebraico e às vezes em aramaico, sem pontuações. Nunca há
página de número 1 em um texto judaico tradicional, os textos sempre
começam na página 2, já que no Talmud não existe nem o começo
nem o fim, tudo é uma questão de meios. O texto principal é posto ao
centro da página rodeado de comentários dos mais diversos estudiosos ---
de diferentes territórios, gerações e contextos. A diagramação é feita
de tal maneira que não há hierarquia entre as interpretações rabínicas,
as exposições dos argumentos e os conceitos apresentados, ou seja, todos
os estudiosos e estudiosas são convidados a trazer uma nova
interpretação das leis judaicas.

Alguns comentadores escreveram tratados, outros apenas alguns rabiscos
sobre uma folha de rascunho, mas no Talmud tudo serve para
tornar as discussões mais dinâmicas. Dessa forma, as mais diversas
visões rabínicas são apresentadas enquanto continuidades dos próprios
textos sagrados, abrangendo pelo menos dois mil anos. Por fim, as
perspectivas rabínicas sobre os textos religiosos e sagrados --- suas
interpretações --- são condicionadas por suas épocas e momentos
históricos, suas relevâncias para a comunidade e seus sentidos judaicos.

\begin{center}\adforn{49}\end{center}

\noindent{}Certa vez, ``perguntaram ao Rabi Mihal: Está escrito na Ética dos
Pais: Quem é sábio? Aquele que aprende com todos os homens, como está
escrito: De todos os meus mestres recebi entendimento. E por que então
não está escrito: Aquele que aprende com todos os mestres?. O Rabi
Mihal explica: O mestre que pronunciou tais palavras pretende explicar
que se deve aprender não só com aqueles que atuam como mestres, mas com
todos os homens. Mesmo com o ignorante, mesmo com o perverso, te é dado
alcançar o entendimento de como conduzir a vida''\footnote{\textsc{buber},
  \emph{op.\,cit.}, p.\,189.}.

Os estudos sobre os livros e textos judaicos religiosos são feitos em
parceria --- \emph{chavruta}\footnote{\emph{Chavruta} é uma palavra
  hebraica proveniente da raíz \emph{chaver}, em português
  ``amigo''. E é uma prática de estudo em que os textos são lidos em
  pares.}, pois se faz necessário falar a respeito do texto como se ele
estivesse realmente vivo. Ao contrário das bibliotecas modernas, lugares
de culto ao silêncio, o \emph{Beit Midrash} (em português: ``Casa de
Estudos'') é o local onde os debates são acalorados, ruidosos, animados
e cheios de gesticulação e de frustrações. A Casa de Estudos reflete o
livro --- o tumulto e o barulho de cada uma das interpretações.

Entende-se aqui que há uma incapacidade de aprender o mundo por completo
em uma só palavra (ou em uma só opinião), pois ``um poema nunca é
concluído, apenas abandonado''\footnote{\textsc{manguel}, \emph{op.\,cit.}, p.\,42.}.
Como observa Flaubert, ``a estupidez consiste em desejar
concluir''\footnote{FLAUBERT \emph{apud} ibidem, p.\,46.}. Ou seja, o texto
será tanto mais significativo quanto maior o número de leituras e de
interpretações. E assim, nem mesmo as diversas interpretações de
\emph{Bereshit} chegam a uma conclusão. Muitas são as traduções e
interpretações de ``No princípio era o verbo'', porém, o mais
correto seria ``Deus começa a criar'', na medida em que o
judaísmo não produz construções que possam levar à idolatria. Rejeita, a
todo custo, a tentação da serpente em aspirar à condição divina, de
fixar e de sacralizar conceitos em uma só conclusão. Se partirmos do
pressuposto de que ``No princípio era o verbo'', então
automaticamente estamos colocando um começo, fixando imagens. Como
resposta todos os livros e, em especial, a máxima literária judaica, a
Torá, são passíveis de interpretações e de apropriações.

% Comentário da editora: De onde é "no princípio era o verbo"? Do próprio Bereshit?

Uma história talmúdica nos conta que, em meio a uma discussão
rabínica, dois indivíduos, buscando defender seus pontos de vista,
iniciam um debate a respeito de uma lei judaica. O primeiro rabino diz:
``Se eu estiver correto que esta árvore o prove''. Depois de
alguns segundos, passa um vento sobre a árvore fazendo-a balançar. O
segundo rabino responde: ``Não há como comprovar que essa árvore
balançou por conta de sua suposição''. O primeiro rabino novamente diz:
``Se eu estiver correto que esta pedra o prove''. Logo em
seguida, a pedra desliza penhasco abaixo em direção ao mar. E novamente
o segundo rabino responde: ``Não há como comprovar que a pedra
deslizou por conta de sua suposição''. O primeiro rabino, impaciente,
diz: ``Se eu estiver correto que os céus o provem''. Nesse
instante a esfera divina brilha, permitindo uma passagem de luz. O
segundo rabino, deslumbrado, observa e diz: ``O poder já não está
mais nas mãos dos céus. A Torá foi dada por Deus aos seres humanos e
agora está sob a vontade deles!''\footnote{\textsc{oz; oz-salzberger}, \emph{Os
  judeus e as palavras}, p.\,18.}.

Durante os últimos dois mil anos, e ainda hoje, as discussões e os
diálogos manifestaram-se através do livro judaico; o debate judaico não
ocorre em fóruns abertos, mas nas páginas escritas. Seus assuntos
abrangem séculos, até mesmo milênios, e convergem de maneiras
excepcionais, criam espaços, territórios de apropriação.

\begin{center}\adforn{49}\end{center}

\noindent{}O \emph{Midrash}\footnote{O \emph{midrash} é uma metodologia judaica de
  leitura das escrituras sagradas. Seu fundamento está em uma visão
  cíclica da história, segundo a qual o acontecimento central, ou
  melhor, a presença de Deus no texto, é continuamente revisado,
  oferecendo sempre novos significados às várias situações expostas, bem
  como às interpretações.} diz que toda e cada letra da Torá deve
estar completamente envolvida por um pedaço de pergaminho\footnote{Referência
  à \emph{Menachot} 29a}. O pergaminho branco em volta das letras pretas
é parte integrante dos escritos e do rolo da Torá. Em debates
talmúdicos extremamente aprofundados, os sábios interpretam que a
Torá é ``escrita em fogo negro, sobre fogo branco''. As
separações e os espaços brancos do pergaminho permitiram que Moisés,
durante a inscrição da fala de Deus no texto, analisasse e absorvesse os
aprendizados anteriores, para assim seguir ao próximo. O fogo branco é o
espaço mais elevado da contemplação, ao passo que o fogo negro (a
escrita) é a revelação do intelecto no domínio da linguagem.

\begin{center}\adforn{49}\end{center}

\noindent{}A cultura e a filosofia judaicas, bem como as suas tradições, conhecem
muito bem a relevância do espaço da página, dos textos e dos livros para
a construção de territorialidades. A divindade judaica habita o texto,
assim como também habita a recordação --- a lembrança\footnote{Por mais
  que este trabalho se sustente em uma perspectiva judaico-diaspórica
  sobre o lugar, ou seja, em territorialidades judaicas móveis,
  transitórias e fluidas --- sendo elas: \emph{BS''D}, \emph{kipá},
  \emph{mezuzá}, \emph{eruv} --- há um território espiritual e simbólico
  que permeia todas as discussões: a memória e a lembrança do Templo
  Sagrado de Jerusalém. Ao final, todas as rezas, orações, tradições e
  cultos judaicos religiosos são conduzidos à cidade de Jerusalém, local
  sacrossanto judaico --- território da construção do Templo ---, e
  assim levados ao céu.} --- e os anseios de retorno do Templo Sagrado
de Jerusalém.

As territorialidades sagradas dos textos e das páginas judaicas
religiosas são cosmogonizadas por meio da inscrição ao canto da folha do
artifício de {בס''ד} (BS''D). Em aramaico, a inscrição BS''D
(acrônimo de: \emph{B'Syata D'bishmaya}) significa ``com a ajuda
dos céus''. Em hebraico moderno, a mesma inscrição significa: ``em uma
cinta''. Este artifício passou a ser introduzido pelo Rabi Yehuda
HaChassid entre os anos de 1150 e 1217, enquanto um símbolo de lembrança
divina. Ou seja, quando realizamos tarefas cotidianas e mundanas Deus
deve estar em nossas mentes, e em nossas penas e canetas\footnote{A
  inscrição da tinta de uma caneta sobre uma página cria um território.}.

O judaísmo entende que as páginas --- estes espaços em branco, prontos
para serem riscados --- são locais de contemplação e de criação. Ao
realizar um traço, um desenho e/\,ou um rascunho sobre um pedaço de papel,
os seres humanos criam um mundo e um universo, por completo --- seja ele
virtual, ficcional ou real. Escrever é criar mundos, e o mesmo pode ser
dito de desenhar, gravar, inserir ou apagar, como nos lembra a história
do \emph{golem}.

Assim, há o costume de gravar o acrônimo BS''D nas páginas,
anteriormente a quaisquer inscrições. Este artifício relembra que apenas
``com a ajuda dos céus'' é possível criar --- ao final, nada está acima
da criação divina, nem mesmo a fruta comida por Adão no Jardim do
\emph{Éden}.

Por fim, ambas as traduções --- ``com a ajuda dos céus'' (do
aramaico ao português) e ``em uma cinta'' (do hebraico moderno ao
português) --- relacionam-se com a Torá. A Torá é um
material escrito que dialoga com a existência, com a crença e com a
presença de Deus (por meio do texto) no mundo humano. Este
livro-documento é envolto e fechado por uma cinta --- um pedaço de pano
ou de couro --- e encapsulado em uma caixa de madeira. Assim, a tradição
judaica se estabelece no enlaçamento (cinta) entre a presença divina, a
fidelidade e a lealdade judaica ao monoteísmo (com a ajuda de Deus) e a
fé no livro.

\begin{center}\adforn{49}\end{center}

\noindent{}O \emph{aleph} ({א}) é a primeira letra do alfabeto hebraico --- e do
alfabeto aramaico ---, usado para escrever o iídiche e o ladino.
É uma letra que contém dentro de si todas as demais letras --- contém o
universo inteiro, sem fim, e ``tem a forma de um homem que
assinala o céu e a terra, para indicar que o mundo inferior é o espelho
e o mapa do superior''\footnote{\textsc{kilsztajn}, \emph{Yiddish}, p.\,20.}. O
\emph{aleph} ({א}) é o Lugar onde estão, sem se confundirem, todos os
lugares do mundo, vistos por todos os ângulos --- onde o interlocutor vê
o que está nascendo e partindo.

Certa vez li que a palavra \emph{vayikra} --- {ויקרא} (em português:
``Ele chamou'') --- está escrita de uma forma pouco comum na
Torá. A última letra, o \emph{aleph}, está em miniatura. Os
sábios interpretam que a tentativa de esconder esta letra foi uma
expressão de extrema humildade e de devoção de Moisés a Deus. Por assim
dizer, Deus apenas ``aconteceu'' (\emph{vayikar}) --- {ויקר}.

% \begin{center}\adforn{49}\end{center}

% \noindent{}O ato de escrever, a incisão feita ao gravar letras, palavras e frases
% sobre uma página em branco, cria mundos, realidades e universos por
% completo. Este incansável e interminável movimento é feito neste exato
% instante e neste exato documento\footnote{Pretende-se que as diversas
%   interpretações dos textos, aqui expostos, ressoem em você que lê, o
%   leitor.}. E assim, como intérprete de uma longa tradição de textos e
% de livros, acredito que o nosso mundo --- e todas as coisas que ele
% contém --- está sendo criado e reinventado por nós a todo instante.

% Por fim, os textos aqui escritos e inscritos --- suas notas de
% rodapé\footnote{A disposição gráfica das notas de rodapé alude à página
%   dos textos judaicos religiosos, bem como do pensamento rabínico que há
%   no Talmud. Pois, no Talmud, a menção de um nome,
%   situação ou ideia permite a introdução de uma história que alivia o
%   clima de discussões complexas e leva o debate adiante. As notas, ou
%   melhor, os comentários e as interpretações rabínicas, rondam,
%   circunscrevem e flutuam ao redor do texto principal sem
%   hierarquização. Aludindo à organização talmúdica, procuro
%   incitar os outros a refletir, abrindo espaço para as mais diversas
%   interpretações, histórias e vozes.} e suas interpretações
% (ressonâncias) que levam a outras histórias, anedotas e contextos ---
% atestam a veracidade da fala de Rabi Tarfon: ``Não lhe é exigido
% que complete a tarefa, mas não é livre para escapar dela''\footnote{\emph{Pirkei
%   Avot} 2:16, escrito cerca do século \textsc{i}.}.

% A diagramação destas páginas-territórios são releituras dos textos
% sagrados, uma tentativa de apropriação das chispas e das fagulhas que
% voam e contornam o território do texto --- já que o espaço do texto é um
% território. E mais que nada, neste instante, é hora de nos apropriarmos
% dele.

\begin{center}\adforn{49}\end{center}

\noindent{}``No mundo da aprendizagem, das discussões e dos debates judaicos,
as palavras dos textos, literal e figurativamente,
ressoam''\footnote{HALPERN, \emph{The Art of Debate: Jewish Style}. Tradução da autora.}. E, ao ressoar, entoam uma voz --- um sopro de
espírito, uma alma. As palavras, finalmente, cosmogonizam o
mundo\footnote{As vozes não se calam, os sopros fazem eco para a
  eternidade.}.

Desde outubro de 2023 \emph{ya no ay biervos}\footnote{Do ladino: ``já
  não há palavras''.}. ``No estupendo Nossa Música, de Jean-Luc
Godard, uma jornalista israelense entrevista o poeta palestino Darwish,
que, privado de sua terra, fez das palavras a sua pátria. E ela comenta:
\emph{você começa a soar como judeu!}. O devir-judeu do palestino, o
devir-palestino do judeu''\footnote{\textsc{pelbart}, Peter Pál. \emph{Ser judeu
  no Brasil}.}.

Por fim, todos nós, antes de mais nada, vivemos em palavras.